% Brüche und Dezimalzahlen · Vergleichen und Ordnen · Prüfungs-Fokus MSA – bank/brueche-dezimalzahlen/e5.jsonl (zusammenbau v0.6, Prüfungs-Fokus)
\einheitenkopf[e5]{Einheit 5 · Vergleichen, Ordnen, Runden}
\zweigzeile{ab Kl. 5, je nach Buch bis Kl. 6 · P10 ×8}
\setcounter{aufgabe}{0}

% Anlauf (ohne Prüfkennung)
\begin{aufgabe}{Vergleichen und Ordnen – Anlauf}
\begin{teile}
\teil Hänge Nullen an, bis beide gleich viele Stellen haben: $0{,}7$ und $0{,}42$. \leerfeld und \leerfeld
\teil Welche Zahl ist größer: $0{,}47$ oder $0{,}52$? \leerfeld
\teil Welche Zahl ist größer: $0{,}8$ oder $0{,}79$? \leerfeld
\end{teile}
\end{aufgabe}

\begin{aufgabe}{Vergleichen und Ordnen – Prüfungsaufgaben}
\begin{teile}
\teil Welcher Wert ist der kleinste: $5{,}5$; $0{,}55$; $0{,}5^2$; $55\,\%$? \hfill (P23) \\ \leerfeld
\teil Welcher Wert ist der kleinste: $0{,}7^2$; $0{,}5$; $4{,}9$; $45\,\%$? \hfill (P23) \\ \leerfeld
\teil Ordne aufsteigend: $-\frac{3}{4}$; $1{,}7$; $-0{,}749$; $\sqrt{3}$ \hfill (P14) \\ \leerfeld $<$ \leerfeld $<$ \leerfeld $<$ \leerfeld
\teil Ordne aufsteigend: $-\frac{2}{5}$; $2{,}2$; $-0{,}41$; $\sqrt{5}$ \hfill (P14) \\ \leerfeld $<$ \leerfeld $<$ \leerfeld $<$ \leerfeld
\teil Welche Aussage ist wahr? Kreuze an. \hfill (P15)\\ \kreuz{$2{,}25 < \frac{9}{4}$}\\ \kreuz{$\frac{12}{5} > \frac{9}{4}$}\\ \kreuz{$\sqrt{5} > \frac{9}{4}$}
\teil Welche Aussage ist wahr? Kreuze an. \hfill (P15)\\ \kreuz{$3{,}25 < \frac{13}{4}$}\\ \kreuz{$\frac{17}{5} > \frac{13}{4}$}\\ \kreuz{$\sqrt{10} > \frac{13}{4}$}
% AUSGELASSEN brueche-dezimalzahlen-e5-k2-s9-v7: Missing $ inserted.
% \teil Setze $<$, $=$ oder $>$ ein: $4\,\%$ __ $0{,}4$ \hfill (P18) \\ \leerfeld
\teil \textit{(ausgelassen)}
% AUSGELASSEN brueche-dezimalzahlen-e5-k2-s9-v8: LaTeX Error: Command \item invalid in math mode.
% \teil Setze $<$, $=$ oder $>$ ein: $0{,}8$ __ $8\,\%$ \hfill (P18) \\ \leerfeld
\teil \textit{(ausgelassen)}
\teil Welche Zahl liegt genau in der Mitte zwischen $-0{,}8$ und $-0{,}7$? \hfill (P20) \\ \leerfeld
\teil Welche Zahl liegt genau in der Mitte zwischen $-1{,}3$ und $-1{,}2$? \hfill (P20) \\ \leerfeld
\end{teile}
\end{aufgabe}

\medskip
\uebersichtskasten{Dezimalzahlen vergleichen: Stelle für Stelle von links – vorher Nullen anhängen, bis alle gleich viele Stellen haben. \quad 0,6 > 0,58 (0,60 > 0,58) \\ Runden: Die Stelle dahinter entscheidet – 0 bis 4 ab, 5 bis 9 auf. \quad 2,749 ≈ 2,7 (Zehntel), ≈ 2,75 (Hundertstel) \\ Mitte zweier Zahlen: beide addieren, durch 2 – oder am Zahlenstrahl eine Stelle feiner einteilen. \quad Mitte von 0,3 und 0,4: 0,35 \\ Verschiedene Darstellungen: alle in Dezimalzahlen umwandeln, dann ordnen. \quad 3/8 = 0,375 \quad 35 \% = 0,35 \quad 0,6² = 0,36 → 0,35 < 0,36 < 0,375}
